% Generated draft: requires mathematical and source review.
\documentclass[11pt]{article}
\usepackage[margin=25mm]{geometry}
\usepackage{fontspec}
\setmainfont{Latin Modern Roman}
\usepackage{amsmath,amssymb,mathtools,unicode-math}
\setmathfont{Latin Modern Math}
\usepackage{enumitem,needspace,xurl,hyperref}
\setlength{\parindent}{0pt}
\setlength{\parskip}{5pt}
\setlength{\emergencystretch}{3em}
\newcommand{\N}{\mathbb{N}}
\newcommand{\Z}{\mathbb{Z}}
\newcommand{\Q}{\mathbb{Q}}
\newcommand{\R}{\mathbb{R}}
\newcommand{\C}{\mathbb{C}}
\newcommand{\F}{\mathbb{F}}
\newcommand{\A}{\mathbb{A}}
\newcommand{\PP}{\mathbb{P}}
\newcommand{\E}{\mathbb{E}}
\newcommand{\eps}{\varepsilon}
\newcommand{\dd}{\,\mathrm{d}}
\begin{document}
\section*{UnsolvedMath ID: 30002289}
\textbf{Source problem code:} catalogue entry
\subsection{Definitions and mathematical statement}
The supplied record concerns local regularity of scalar \(p\)-harmonic functions. It does not specify the range of \(p\) or the solution class. The following is therefore a conditional formulation using the usual conventions \(1<p<\infty\) and real-valued weak solutions; these conventions require confirmation against the cited report.

Let \(n\geq 3\), let \(\Omega\subset\R^n\) be open, and let \(u\in W^{1,p}_{\mathrm{loc}}(\Omega)\) satisfy
\[
\int_\Omega |\nabla u|^{p-2}\nabla u\cdot\nabla\varphi\,\dd x=0
\qquad\text{for every }\varphi\in C_c^\infty(\Omega).
\]
Here \(|\xi|^{p-2}\xi\) is assigned the value zero at \(\xi=0\). Define
\[
V_p(\xi)=
\begin{cases}
|\xi|^{(p-2)/2}\xi,&\xi\neq0,\\
0,&\xi=0.
\end{cases}
\]
For \(0<\alpha\leq1\), membership in \(C^{0,\alpha}_{\mathrm{loc}}(\Omega;\R^n)\) means that there is a representative \(f\) such that, for every compact \(K\subset\Omega\), some finite constant \(C_K\) satisfies
\[
|f(x)-f(y)|\leq C_K|x-y|^\alpha
\qquad\text{for all }x,y\in K.
\]

Under these conventions, the explicit conjectural question is whether every such \(p\)-harmonic function satisfies
\[
V_p(\nabla u)\in C^{0,1/2}_{\mathrm{loc}}(\Omega;\R^n)
\]
for every \(n\geq3\) and every admissible \(p\). The source also describes this conjecture as holding ``in every dimension,'' without explicitly specifying its dimension range. No boundary regularity is requested, and no boundary conditions are supplied.

The accompanying problem asks for higher-dimensional analogues of the planar least-regular examples and for the optimal local H\"older regularity of \(\nabla u\). A possible precise interpretation is to determine, for each fixed \((n,p)\), the supremum of exponents \(\alpha\in(0,1]\) for which every such solution has \(\nabla u\in C^{0,\alpha}_{\mathrm{loc}}\), determine whether that supremum is attained, and construct examples establishing sharpness. The supplied record does not describe the planar construction or specify what properties its higher-dimensional analogues must retain, so this interpretation is provisional.

The original statement additionally asks for endpoint and finer regularity of \(p\)-Laplacian solutions as \(q\to\infty\). No definition of \(q\), associated equation, data class, or estimates is supplied. That additional target cannot be formulated precisely from this record and is not identified here with the corrected \(p\)-harmonic formulation.
\subsection{Short English statement}
For higher-dimensional \(p\)-harmonic functions, determine the sharp local H\"older regularity of the gradient and construct examples showing sharpness. The explicit conjecture asks whether \(|\nabla u|^{(p-2)/2}\nabla u\) always has a locally \(1/2\)-H\"older continuous representative. The range of \(p\), intended planar examples, and the original \(q\to\infty\) question need clarification. This is a statement-only draft requiring human mathematical review.
\subsection{Sources}
\begin{itemize}
\item Local UnsolvedMath registry, permanent ID \(30002289\), source problem code \texttt{catalogue entry}, title: Endpoint Regularity for \(p\)-Laplacian Gradients. The supplied record belongs to the Oberwolfach Reports Open Problems collection and lists the proposed year as 2013; no individual proposer is supplied.
\item Supplied report citation: \emph{Mini-Workshop: The p-Laplacian Operator and Applications} (2013), Oberwolfach Reports, pp.\ 433--482, DOI \texttt{10.4171/owr/2013/08}. Source URL and supplied literature-source link, labeled ``Oberwolfach Report 2013'': \url{https://doi.org/10.4171/owr/2013/08}.
\item The record labels the statement \texttt{corrected\_verified} and reports that the official open-problem session asks for higher-dimensional analogues of the planar construction and explicitly conjectures local \(C^{1/2}\) regularity of \(V_p(\nabla u)\). This is supplied verification metadata; the report has not been independently consulted for this draft.
\item Literature assessment dated exactly \texttt{2026-08-21}: the supplied status is ``open,'' and the supplied assessment says that the planar optimal exponent is known while the sharp gradient H\"older exponent and matching higher-dimensional extremals for \(n\geq3\) remain unresolved. This is supplied dated curation derived from the OWR record, not an independently verified literature assessment. Its time-sensitive status and cited claims require independent verification.
\end{itemize}
\end{document}
